\documentclass[11pt]{article}
\usepackage[T1]{fontenc}
\usepackage[utf8]{inputenc}
\usepackage{lmodern,amsmath,amssymb,amsthm,booktabs,tabularx,microtype}
\usepackage[a4paper,margin=27mm]{geometry}
\usepackage[hidelinks]{hyperref}
\usepackage{fancyhdr}
\pagestyle{fancy}\fancyhf{}\fancyhead[L]{\small Erd\H{o}s 727: a proposed $k=3$ extension}\fancyhead[R]{\small 17 September 2026}\fancyfoot[C]{\thepage}
\setlength{\headheight}{14pt}\setlength{\parindent}{0pt}\setlength{\parskip}{6pt}
\newtheorem{theorem}{Proposed theorem}
\newtheorem{lemma}{Lemma}
\newcommand{\ind}{\mathbf1}\newcommand{\MX}{\mathcal M_X}\newcommand{\NX}{N_X}
\title{Six linear factors and a positive parameter bound:\\a proposed extension to $k=3$ in Erd\H{o}s problem 727}
\author{AI-assisted mathematical continuation}
\date{17 September 2026}
\begin{document}
\maketitle
\begin{abstract}
The supplied manuscript treats only $k=2$. This separate continuation uses
$210=14\cdot15$ and $420=20\cdot21$ to factor all three consecutive integers
$n+1,n+2,n+3$ for the family $n=210t^2+391t+179$. Six simple root classes replace
the earlier four. A signed residue calculation at modulus 210 gives an average
large-prime obstruction coefficient $5/6$. Combining this with uniform carry
estimates gives the proposed lower limiting success proportion
$0.255039321662196649\ldots$ for $((n+3)!)^2\mid(2n)!$. The argument is written
below with its changed hypotheses and costs, not inferred from finite scans.
\end{abstract}
\textbf{Status and scope.} This is a new proposed $k=3$ argument developed during
this review, not a result asserted in the supplied September 16 PDF or the older
repository. It has not received independent peer review or formal verification.
No claim of priority, accepted resolution, or a result for every fixed $k$ is
made. Numerical checks in Section~\ref{computations} verify finite arithmetic
only. All constants, including modulus 210, are fixed before the limiting
parameter tends to infinity.

\section{The explicit family and proposed bound}
Put
\begin{equation}\label{family}
 H(t)=210t^2+391t+179,\qquad m=210t+195,\qquad n=H(t)=\frac{m(m+1)}{210}-3.
\end{equation}
As in the $k=2$ companion, write
\begin{equation}\label{S}
 S=4\sum_{J=4}^{\infty}2^{1-J}\log\frac{J+1}{J}.
\end{equation}
All logarithms are natural. Define
\begin{equation}\label{rho}
 \rho_3=1-\frac32 S-\frac5{12}\log3
       =0.2550393216621966492464063134168790975\ldots .
\end{equation}
\begin{theorem}\label{main}
The explicit family \eqref{family} satisfies
\begin{equation}\label{bound}
 \liminf_{T\to\infty}\frac1T
 \#\{t\in[T,2T]\cap\mathbb Z:((H(t)+3)!)^2\mid(2H(t))!\}\ge\rho_3>0.
\end{equation}
Consequently the stated argument, if correct, yields infinitely many distinct
positive solutions for $k=3$.
\end{theorem}
This statement does not upgrade the parameter bound for the original polynomial
$6t^2+11t+3$, nor does it assert any finite interval must attain $\rho_3$.

\section{All three denominator factors are controlled}
The choice of 210 permits the simultaneous identities
\begin{align}\label{sixfactors}
 n+1&=\frac{(m-20)(m+21)}{210}=(6t+5)(35t+36),\nonumber\\
 n+2&=\frac{(m-14)(m+15)}{210}=(t+1)(210t+181),\\
 n+3&=\frac{m(m+1)}{210}=(14t+13)(15t+14).\nonumber
\end{align}
Each of the six displayed linear polynomials is primitive. Thus the denominator
product $E(n)=(n+1)(n+2)(n+3)$ has exactly these six primitive factors, including
multiplicities. Cancelling $(n!)^2$ gives
\begin{equation}\label{criterion}
 ((n+3)!)^2\mid(2n)!\iff E(n)^2\mid\binom{2n}{n}.
\end{equation}
For a prime $p$, Legendre's formula gives
\begin{equation}\label{carries}
 C_p(n):=v_p\binom{2n}{n}
 =\sum_{j\ge1}\ind_{\{\{n/p^j\}\ge1/2\}}.
\end{equation}
It suffices to check $C_p(n)\ge2v_p(E(n))$ at primes dividing $E(n)$.

Let
\[
 \mathcal C=\{0,1,-14,15,-20,21\},\qquad d_c=1-2c,
\]
\[
 i_c=\begin{cases}3,&c=0,1,\\2,&c=-14,15,\\1,&c=-20,21.\end{cases}
\]
The maximum difference between two shifts is 41. For $p>41$, the prime is a unit
modulo 210 and divides at most one $m+c$. If $p^e\mid m+c$, then
$v_p(E(n))=e$ and $n\equiv-i_c\pmod{p^e}$. Since $i_c\le3$ and $p>41$,
\[
 \{n/p^j\}=1-i_c/p^j>1/2\quad(1\le j\le e).
\]
There are at least $e$ initial carries. In particular, when $v_p(E(n))=1$,
failure is equivalent to $C_p(n)=1$: every later carry must be absent.

For $m+c=ap$, direct substitution gives
\begin{equation}\label{ap}
 n=\frac{a^2p^2+d_cap}{210}-i_c,\qquad
 \frac n{p^2}=\frac{a^2}{210}+\frac{d_ca}{210p}-\frac{i_c}{p^2}.
\end{equation}
The six derivative constants are $1,-1,29,-29,41,-41$, all nonzero modulo
primes $p>41$.

\section{Fixed primes, repeated primes and small growing primes}
Work on
\[
 \MX=\{m\in[X,2X]\cap\mathbb Z:m\equiv195\pmod{210}\},\qquad
 \NX=X/210+O(1).
\]
\subsection{Every fixed prime has a density-zero failure set}
For each fixed $p$, primitivity in \eqref{sixfactors} implies
\[
 \limsup_{T\to\infty}\frac1T\#\{t\in[T,2T]:v_p(E(H(t)))\ge6h\}\le6p^{-h}.
\]
Each factor has either no root or one root modulo $p^h$. A product valuation
at least $6h$ requires a factor valuation at least $h$.

For fixed $K$, the number of residues modulo $p^L$ with at most $K$ carries in
the first $L$ levels is
\[
 O_{p,K}\bigl(L^K\lceil p/2\rceil^L\bigr)=o(p^L).
\]
Indeed a digit at least $\lceil p/2\rceil$ forces a carry whatever the incoming
carry; choose at most $K$ such digit positions. For $p\mid210$,
$H(t)\equiv391t+179\pmod p$ is a permutation and
$H'(t)=420t+391$ is always a unit. Simple lifting makes $H$ a permutation
modulo every $p^L$.

For $p\nmid210$, the prime is odd and the exact identity
\begin{equation}\label{square}
 840H(t)+2521=(420t+391)^2
\end{equation}
reduces the issue to square-root fibers. Discard $p^h\mid420t+391$, a class of
density $p^{-h}$. For $L>2h$ the remaining fibers have at most $2p^h$ elements:
a root with valuation $s<h$ reduces to a unit square-root congruence modulo
$p^{L-2s}$, with at most two roots and $p^s$ lifts. Send $L\to\infty$ first,
then $h\to\infty$. This proves $C_p(H(t))\to\infty$ in probability. Combining
it with the denominator valuation bound proves the fixed-prime assertion.
All primes at most 41 are included here; no simplicity assertion is needed
for those primes.

\subsection{Repeated nonfixed factors}
Fix $P>41$. For $p>P$, $p^2\mid m+c$ specifies one class modulo $210p^2$ and
requires $p\le\sqrt{2X+21}$. Summing the progression count
$X/(210p^2)+O(1)$ over six shifts and these primes gives
\[
 O(X/P)+O(\sqrt X).
\]
Its normalized limsup is $O(1/P)$. Outside this exception and the fixed-prime
exceptions, all obstructing primes have $v_p(E(n))=1$.

\subsection{Small growing primes}
For $P<p\le X^\delta$, $0<\delta<1/4$, the six root classes in the $t$ variable
are distinct and simple, with derivatives $d_c$ modulo $p$. On one root class,
$H(t)$ maps bijectively modulo $p^L$ onto $n\equiv-i_c\pmod p$. Take
$L=\lfloor\log_p(X/420)\rfloor$. The first carry is forced and each later
no-carry digit has at most $(p+1)/2$ choices. Each permissible residue appears
$O(X/p^L)$ times in the parameter interval. Since $(p+1)/(2p)\le4/7$, the
bad-incidence count per root is at most
\[
 \frac{A X}p\exp\left(-c_0\frac{\log X}{\log p}\right),\qquad c_0=\log(7/4),
\]
for a fixed absolute constant $A$ (210 is fixed). The partial summation in
Lemma 6.1 of the supplied draft\cite{draft} bounds the normalized total over
six roots by a fixed constant times
\[
 \int_{1/\delta}^{\infty}e^{-c_0v}\frac{dv}{v}=: \varepsilon(\delta)/A',
\]
up to a term tending to zero with $X$. In particular
$\varepsilon(\delta)\to0$ as $\delta\downarrow0$. The argument uses root
lifting, not independence among the six factors.

\section{Uniform joint carry estimates}
For a fixed prime and shift, $ap-c\equiv195\pmod{210}$ implies
$a=a_0+210b$. The variable $b$ runs through an interval of
\begin{equation}\label{length}
 L_{p,c}=X/(210p)+O(1)
\end{equation}
consecutive integers. Fix $\delta,\eta>0$, $J\ge4$, and a nonzero integer
frequency vector $h=(h_2,\ldots,h_J)$. Suppose
\begin{equation}\label{range}
 X^\delta\le p\le X^{1/2-\eta},\qquad p\le X^{2/J-\eta}.
\end{equation}
In the phase $\sum_{j=2}^J h_j n/p^j$, remove the integer-valued variation of
$h_2a^2/210$. If all higher frequencies vanish, the phase is linear with slope
$h_2d_c/p$; its normalized geometric sum is $O_h(p^2/X)=o(1)$.

Otherwise, with $j_0\ge3$ the smallest nonzero higher frequency, the remaining
second derivative is
\[
 420\sum_{j=3}^J h_jp^{2-j}
 =420h_{j_0}p^{2-j_0}(1+O_h(p^{-1})).
\]
The first nonzero term dominates for fixed $h$ and growing $p$. The standard
second-derivative estimate\cite{robert}, divided by $L_{p,c}\asymp X/p$, gives
\[
 O_h\left(p^{-(j_0-2)/2}+p^{j_0/2}/X\right)
 =O_h\left(X^{-\delta/2}+X^{-J\eta/2}\right)=o(1).
\]
These estimates are uniform in the primes in \eqref{range} and all six shifts.
The fixed finite trigonometric majorant/minorant argument of draft Lemma
7.1 converts them into uniform box counts. Thus
\begin{equation}\label{box}
 \frac1{L_{p,c}}\#\{a:\{n/p^j\}<1/2\ (2\le j\le J)\}
 =2^{1-J}+o_{\delta,\eta,J}(1).
\end{equation}
No independence of digit levels at a fixed prime is asserted. The Fourier
frequencies are fixed first, so the exact relations with frequencies growing
with $p$ do not contradict this argument.

For $X^{1/2+\eta}\le p\le X^{2/3-\eta}$, the single phase $h n/p^3$ has second
derivative $420h/p$. Its normalized exponential sum is
$O_h(p^{-1/2}+p^{3/2}/X)=o(1)$, uniformly in all six shifts. It follows that
\begin{equation}\label{middle}
 \frac1{L_{p,c}}\#\{a:\{n/p^3\}<1/2\}=1/2+o_\eta(1)
\end{equation}
in this intermediate range.

\section{The signed residue calculation at large primes}
If $p\ge X^{1/2+\eta}$, then $a/p=O(X^{-2\eta})$. In \eqref{ap}, the leading
fractional part is $s/210$, where
\[
 s\equiv(195+c)^2p^{-2}\pmod{210},\qquad 0\le s<210.
\]
At $s=0$ the fraction is below one half exactly when $d_c>0$. At $s=105$ it
is below one half exactly when $d_c<0$. Away from these boundaries it is below
one half exactly when $0<s<105$. These signs are uniform for large $X$ because
$d_cap-210i_c$ has the sign of $d_c$ as $ap=m+c\to\infty$.

The 48 units modulo 210 have six square classes, each of multiplicity eight.
The table lists the numerator $s$ for each of the six shifts; its final column
counts the shifts which fail to supply the second carry.
\begin{center}
\small
\begin{tabular}{r|rrrrrr|r}
\toprule
$p^{-2}\bmod210$ & $c=0$ & $1$ & $-14$ & $15$ & $-20$ & $21$ & bad count\\
\midrule
1   &15 &196&1  &0&175&36 &3\\
79  &135&154&79 &0&175&114&1\\
109 &165&154&109&0&175&144&0\\
121 &135&196&121&0&175&156&0\\
151 &165&196&151&0&175&186&0\\
169 &15 &154&169&0&175&204&1\\
\bottomrule
\end{tabular}
\end{center}
In particular the zero column $c=15$ is \emph{not} bad: its negative correction
puts the fraction just below one, not near zero from above. The average bad
count over prime unit classes is the exact rational value
\begin{equation}\label{b}
 b_3=\frac{8(3+1+0+0+0+1)}{48}=\frac56.
\end{equation}
This table is derived from exact modular arithmetic, not fitted to scan data.

\section{Prime summation, boundary coverage and completion}
For fixed unit $v\pmod{210}$ and $0<\alpha<\beta$, the unconditional
fixed-modulus reciprocal-prime asymptotic\cite[eq.~(6)]{keliher} gives
\begin{equation}\label{AP}
 \sum_{\substack{X^\alpha<p\le X^\beta\\p\equiv v\; (210)}}\frac1p
 =\frac1{48}\log\frac\beta\alpha+o(1).
\end{equation}
This is the classical fixed-modulus formula, not the GRH-dependent varying-modulus
estimates later in that reference. An upper endpoint $2X+21$ is allowed with
$\beta=1$. By \eqref{length}, the total $O(1)$ progression errors contribute
$O(\pi(2X+21)/X)=o(1)$ after normalization. On primes $p\ge X^\delta$,
a uniform conditional error $o(1)$ sums to $o(1)$ because the reciprocal-prime
sum is bounded at fixed $\delta$.

Below $X^{1/2}$, omit narrow exponent-boundary neighborhoods and use the bands
\[
 X^{\max\{\delta,2/(J+1)\}+\eta}<p<X^{2/J-\eta},\qquad
 J\ge4,\quad 2/J>\delta.
\]
There are finitely many at fixed $\delta$. A bad valuation-one prime must have
no carry at levels $2,\ldots,J$, so \eqref{box} bounds its cost per shift by
$2^{1-J}\log((J+1)/J)+o(1)$. Summing six shifts and then letting the cutoff
shrink gives at most
\[
 6\sum_{J\ge4}2^{1-J}\log\frac{J+1}{J}=\frac32 S.
\]
For the intermediate primes of \eqref{middle}, only the residue-selected bad
shifts can obstruct. Equation \eqref{AP}, the table, and \eqref{middle} bound
their cost by
\[
 \frac{b_3}{2}\log\frac{2/3-\eta}{1/2+\eta}+o(1).
\]
For $p>X^{2/3-\eta}$, count every incidence in a residue-selected bad shift,
without imposing a further carry condition. The cost is at most
\[
 b_3\log\frac1{2/3-\eta}+o(1).
\]
As $\eta\downarrow0$ these two costs sum to
$b_3(\log(3/2)+\tfrac12\log(4/3))=(5/12)\log3$.

For explicit coverage, let
\[
 \mathcal E_\delta=\{\delta,1/2,2/3\}\cup
 \{2/J:J\ge4,\ \delta\le2/J\le1/2\}.
\]
For $0<\eta<\min(\delta/2,1/12)$ the union of the bands
$[X^{\alpha-\eta},X^{\alpha+\eta}]$, $\alpha\in\mathcal E_\delta$, has total
normalized incidence cost at most
\[
 6\sum_{\alpha\in\mathcal E_\delta}
 \log\frac{\alpha+\eta}{\alpha-\eta}+o(1).
\]
This tends to zero as $\eta\downarrow0$ at fixed $\delta$. The lower bands
partition the remaining exponents between $\delta$ and $1/2$; the intermediate
and large ranges cover the rest. There is no uncovered exponent interval.

Take $X\to\infty$ first with $P,\delta,\eta$ fixed. The failure limsup is at
most the sum of the fixed/repeated cost $O(1/P)$, the small-prime cost
$\varepsilon(\delta)$, the boundary cost just displayed, the finite lower-band
sum, and the intermediate/large-prime costs. Next take $\eta\downarrow0$, then
$\delta\downarrow0$, then $P\to\infty$. The resulting upper bound for failure
is $(3/2)S+(5/12)\log3=1-\rho_3$.

Finally $X=210T$ gives the $t$ interval
$[T-195/210,2T-195/210]$, differing from $[T,2T]$ by $O(1)$ integers, and
$\NX=T+O(1)$. Taking complements proves \eqref{bound}. Strict positivity needs
no decimal evaluation: $S\le1/4$ and $\log3<11/10$ give
\[
 \rho_3\ge\frac58-\frac5{12}\log3>\frac16>0.
\]
The polynomial $H$ is strictly increasing for $t\ge0$, so the successful
parameters give infinitely many distinct $n$. This concludes the proposed proof.

\section{Finite checks and remaining review}\label{computations}
Two implementations use different factorization and carry computations:
Python trial division with explicit digit addition, and a C++ smallest-prime-factor
sieve with Legendre sums. They agree on the first 1,000 parameters and on a
separate fixed 1,000-parameter high-range comparison. The larger C++ scans give
\begin{center}
\begin{tabular}{rrr}
\toprule
Parameter interval & valid parameters & observed fraction\\
\midrule
$[0,1000)$ &126&0.12600\\
$[10000,20000)$ &1880&0.18800\\
$[100000,200000)$ &21456&0.21456\\
\bottomrule
\end{tabular}
\end{center}
The cross-check on $[199000,200000)$ agrees in all 1,000 cases, of which 240
are valid. At $t=61$, $m=13005$ and $n=805440$; both methods certify
\[
 (805443!)^2\mid1610880!.
\]
The prime-by-prime certificate is saved separately. A literal-factorial check
for every $0\le n\le500$ agrees with the binomial criterion and carry test.
These statements are finite arithmetic evidence, not a numerical proof of the
Fourier estimates, prime sums, or limit theorem.

The highest-priority independent checks are the six-factor adaptation of root
lifting, the signed large-prime table (especially its zero column), fixed-modulus
prime weighting, and the uniform-error and limit-order arguments. The next
unresolved extension is not obtained by merely incrementing $k$: one would need
to control $n+4$ as well. This document makes no assertion for $k\ge4$.

\begin{thebibliography}{9}
\bibitem{draft} Supplied candidate manuscript, \emph{A quadratic-family approach
to the $k=2$ case of Erd\H{o}s problem 727}, September 16, 2026, 12 pages.
Preserved at \path{papers/quadratic-family-expanded-2026-09-16.pdf}.
\bibitem{general} Companion continuation, \emph{A stronger parameter bound in a
different family: a generalized continuation for Erd\H{o}s problem 727, $k=2$},
September 17, 2026. Preserved at
\path{papers/generalized-families-2026-09-17.pdf}.
\bibitem{robert} O. Robert, \emph{On van der Corput's $k$-th derivative test for
exponential sums}, Indagationes Mathematicae 27 (2016), 559--589,
Theorem 1. DOI: 10.1016/j.indag.2015.11.009.
\bibitem{keliher} D. Keliher and E. S. Lee, \emph{On the constants in Mertens'
theorems for primes in arithmetic progressions}, arXiv:2306.09981v2 (2024),
equation (6), page 2. The unconditional fixed-modulus formula only is used.
\end{thebibliography}
\end{document}
